\scanpage{33}\mspage{VII}
\uline{Lemme} VII : Si $\Gamma$ est un J-polygone on a $\operatorname{ext}\partial E(\Gamma)\subset\operatorname{ext}\Gamma$.

\uline{Démonstration} : Comme $E(\Gamma)$ est un ensemble convexe $\operatorname{ext}\partial E(\Gamma)=\mathbb{R}^{2}-E(\Gamma)$ (d'aprés la preuve du lemme VI) donc $\operatorname{ext}\partial E(\Gamma)\cap\Gamma=\varnothing$ puisque $\Gamma\subset E(\Gamma)$.

Il est clair que $\operatorname{ext}\Gamma\cap\operatorname{ext}\partial E(\Gamma)\ne\varnothing$ (par exemple en prenant un triangle qui contient $R_{\Gamma}$ et $R_{\partial E(\Gamma)}$ à l'interieur on voit que l'exterieur est \unsure{comun|commun} à \deleted{\unsure{à|la}} $\operatorname{ext}\Gamma$ et a $\operatorname{ext}\partial E(\Gamma)$). Donc $\operatorname{ext}\partial E(\Gamma)\subset\operatorname{ext}\Gamma$ puisque $\operatorname{ext}\partial E(\Gamma)$ est connexe \inserted{et $\mathbb{R}^{2}-\Gamma=\operatorname{int}\Gamma\cup\operatorname{ext}\Gamma$.} cela prouve le lemme.

% slanted
\margin{\unsure{inutile|inutil}}
\uline{Lemme VII} : Si $\Gamma$ est un \inserted{(}J-\inserted{)}polygone non-convexe il existe $s,s'\in\mathbb{S}(\Gamma)$ tel que \deleted{\unsure{si}} $s'$ est visible de $s$ de l'exterieur de $\Gamma$.

\uline{Demonstration} : On sait que $\partial E(\Gamma)\cap\Gamma\ne\varnothing$ (puisque $\deleted{\unsure{\exists}}\partial E(\Gamma)\cap\Gamma\supset\mathbb{S}(\partial E(\Gamma))$) on veut montrer \inserted{que} $\partial E(\Gamma)\not\subset\Gamma$. Si $\partial E(\Gamma)\subset\Gamma$ alors \inserted{on prouve facilement que} $\partial E(\Gamma)$ est un sous-ensemble ouvert de $\Gamma$ et en plus fermé puisqu'il est compact donc $\partial E(\Gamma)=\Gamma$ en raison de la conexité de $\Gamma$, ce qui contredit le fait que $\Gamma$ est non-convexe.

% margin note written sideways
\margin{\deleted{\unsure{Compl.|Conclure|Concl.} \unsure{que} $\partial E(\Gamma)\subset\unsure{\Gamma}\ \unsure{\Rightarrow}\ \underset{\substack{\Vert\\ E(\Gamma)}}{E(\partial E(\Gamma))}\subset$ \unsure{bref}}}
Soit $\alpha$\inserted{$=[s,s']$} une arête de $\partial E(\Gamma)$ telleque $\alpha\not\subset\Gamma$. Si $]s,s'[\subset\mathbb{R}^{2}-\Gamma$ cela prouve le lemme puisque $s,s'\in\mathbb{S}(\Gamma)$.

Si $s$ est un point isolé de $\alpha\cap\Gamma$ Soit $t$ le point de $(\Gamma-s)\cap\alpha$ qui réalise la distance de $s$ a l'ensemble compact $\alpha\cap(\Gamma-s)=\alpha\cap\Gamma-\underbrace{\{s\}}_{\inserted{\text{ouvert dans }\alpha\cap\Gamma}}$. Il est clair par le choix de $t$ que $]s,t[\subset\mathbb{R}^{2}-\Gamma$. Prouvons que $t\in\mathbb{S}(\Gamma)$. Si $t\notin\mathbb{S}(\Gamma)$ il existe $\alpha'\in A(\Gamma)$ tel que $t\in$ au segment ouvert correspondant $=\mathring{\alpha}'$, mais alors comme $t\in\mathring{\alpha}'\cap\mathring{\alpha}$ on doit pouvoir trouver un ouvert $\subset\mathring{\alpha}'\cap\mathring{\alpha}$ \inserted{auquel $t$ appartient} cela contredit le choix de $t$, puisque $\mathring{\alpha}'\cap\mathring{\alpha}\subset\alpha\cap\Gamma$. Si $s$ n'est pas isolé on peut trouver une arête $[s,t_1]\in\unsure{A.|A}(\Gamma)$ telle que $[s,t_1]\subset\alpha$. Si $t_1$ est isolé on \unsure{réfait|réfoit} l'argument ci-dessus, dans le cas contraire on considère $[t_1,t_2]$-{}-etc. jusqu'à trouver $t$ isolé dans $[s,s']$.

Montrons maintenant que si $s,s'\in S(\Gamma)$ et $]s,s'[\subset\mathbb{R}^{2}-\Gamma$ si $[s,s']\subset\alpha\subset\partial E\unsure{[|(}\Gamma\unsure{]|)}$ $]s,s'[\subset\operatorname{ext}\Gamma$.
% archive sheet number: 32
